% GENERATED FILE. Edit canonical lessons, then run npm run generate:books.
% canonical-source-hash: fnv1a64:414d15cbf143c092
% canonical-lessons: SA-C326-dipavalih, SA-C326-dipayati, SA-C326-mandayati, SA-C326-kautukam, SA-C326-sobha

\chapter{A Night of Lights}
\label{ch:sa-a-night-of-lights}
\hlchaptermodality{326}

\begin{chapteropening}
\textbf{By the end of this chapter:} \emph{I can describe a festival night of lights.}

Complete the last lesson of chapter 326: I can describe a festival night of lights.
\end{chapteropening}

\section[\texorpdfstring{\sk{दीपावलिः} (dīpāvaliḥ)}{दीपावलिः (dīpāvaliḥ)}]{\sk{दीपावलिः} (dīpāvaliḥ) — the festival of lamps; a row of lamps}

\label{lesson:SA-C326-dipavalih}



\begin{quote}
Before the new one: say the Sanskrit for a charioteer, then the Sanskrit for pierces.
\end{quote}

\subsection*{You'll want to know: \sk{दीपावलिः}}
\textbf{\sk{दीपावलिः}} — \emph{dīpāvaliḥ} — \textquotedblleft{}the festival of lamps; a row of lamps\textquotedblright{}.

\begin{quote}
\sk{दीपावल्यां} \sk{सर्वत्र} \sk{दीपाः} \sk{ज्वलन्ति।}
\end{quote}

\begin{quote}
\emph{dīpāvalyāṃ sarvatra dīpāḥ jvalanti.}
\end{quote}

\begin{quote}
At Diwali lamps burn everywhere.
\end{quote}

A feminine noun.

\begin{cousinweb}
\textbf{\sk{दीपावलिः}} (\emph{dīpāvaliḥ}) is built from \emph{dīpa}, \textquotedblleft{}a lamp\textquotedblright{}, and \emph{āvali}, \textquotedblleft{}a row\textquotedblright{}; MW gives \emph{dīpāvali} as a row of lights, a night illumination, and the festival of Diwali is its modern use (MW).
\end{cousinweb}

\begin{grammarlens}[title={What you've built}]
One more word for telling what happened.
\end{grammarlens}

\subsection*{Your turn}
\begin{itemize}
\raggedright
  \item \emph{Say it:} \emph{dīpāvaliḥ}
  \item \emph{Say it:} \emph{dīpāvaliḥ}, once more
  \item \emph{Say it:} \emph{dīpāvalyāṃ sarvatra dīpāḥ jvalanti.} — At Diwali lamps burn everywhere.
  \item \emph{Recall:} say the Sanskrit for a charioteer, then the Sanskrit for pierces, then say \emph{dīpāvaliḥ} again
\end{itemize}

\begin{tcolorbox}[breakable,colback=teal!4,colframe=teal!35!black,title={Before you move on}]
What does \sk{दीपावलिः} mean? (\textquotedblleft{}The festival of lamps; a row of lamps\textquotedblright{}.) Say it once more.
\end{tcolorbox}
\section[\texorpdfstring{\sk{दीपयति} (dīpayati)}{दीपयति (dīpayati)}]{\sk{दीपयति} (dīpayati) — lights; kindles}

\label{lesson:SA-C326-dipayati}



\begin{quote}
Before the new one: say the Sanskrit for pierces, then the Sanskrit for the festival of lamps.
\end{quote}

\subsection*{You'll want to know: \sk{दीपयति}}
\textbf{\sk{दीपयति}} — \emph{dīpayati} — \textquotedblleft{}lights; kindles\textquotedblright{}.

\begin{quote}
\sk{बालिका} \sk{गृहस्य} \sk{द्वारे} \sk{दीपं} \sk{दीपयति।}
\end{quote}

\begin{quote}
\emph{bālikā gṛhasya dvāre dīpaṃ dīpayati.}
\end{quote}

\begin{quote}
The girl lights a lamp at the door of the house.
\end{quote}

Past (she …): \sk{सा} \sk{दीपं} \sk{दीपितवती}, \emph{sā dīpaṃ dīpitavatī}.

\begin{cousinweb}
\textbf{\sk{दीपयति}} (\emph{dīpayati}) comes from the root √dīp, \textquotedblleft{}to blaze, shine\textquotedblright{}, as in \emph{dīpa}, \textquotedblleft{}a lamp\textquotedblright{}, in its causative form: to set alight, kindle (MW).
\end{cousinweb}

\begin{grammarlens}[title={What you've built}]
One more word for telling what happened.
\end{grammarlens}

\subsection*{Your turn}
\begin{itemize}
\raggedright
  \item \emph{Say it:} \emph{dīpayati}
  \item \emph{Say it:} \emph{dīpayati}, once more
  \item \emph{Say it:} \emph{bālikā gṛhasya dvāre dīpaṃ dīpayati.} — The girl lights a lamp at the door of the house.
  \item \emph{Recall:} say the Sanskrit for pierces, then the Sanskrit for the festival of lamps, then say \emph{dīpayati} again
\end{itemize}

\begin{tcolorbox}[breakable,colback=teal!4,colframe=teal!35!black,title={Before you move on}]
What does \sk{दीपयति} mean? (\textquotedblleft{}Lights; kindles\textquotedblright{}.) Say it once more.
\end{tcolorbox}
\section[\texorpdfstring{\sk{मण्डयति} (maṇḍayati)}{मण्डयति (maṇḍayati)}]{\sk{मण्डयति} (maṇḍayati) — decorates; adorns}

\label{lesson:SA-C326-mandayati}



\begin{quote}
Before the new one: say the Sanskrit for the festival of lamps, then the Sanskrit for lights.
\end{quote}

\subsection*{You'll want to know: \sk{मण्डयति}}
\textbf{\sk{मण्डयति}} — \emph{maṇḍayati} — \textquotedblleft{}decorates; adorns\textquotedblright{}.

\begin{quote}
\sk{माता} \sk{गृहं} \sk{पुष्पैः} \sk{मण्डयति।}
\end{quote}

\begin{quote}
\emph{mātā gṛhaṃ puṣpaiḥ maṇḍayati.}
\end{quote}

\begin{quote}
Mother decorates the house with flowers.
\end{quote}

Past (she …): \sk{सा} \sk{गृहं} \sk{मण्डितवती}, \emph{sā gṛhaṃ maṇḍitavatī}.

\begin{cousinweb}
\textbf{\sk{मण्डयति}} (\emph{maṇḍayati}) comes from the root √maṇḍ, \textquotedblleft{}to deck, adorn\textquotedblright{}; MW gives its causative \emph{maṇḍayati} as to adorn, and \emph{maṇḍita} as adorned, decorated (MW).
\end{cousinweb}

\begin{grammarlens}[title={What you've built}]
One more word for telling what happened.
\end{grammarlens}

\subsection*{Your turn}
\begin{itemize}
\raggedright
  \item \emph{Say it:} \emph{maṇḍayati}
  \item \emph{Say it:} \emph{maṇḍayati}, once more
  \item \emph{Say it:} \emph{mātā gṛhaṃ puṣpaiḥ maṇḍayati.} — Mother decorates the house with flowers.
  \item \emph{Recall:} say the Sanskrit for the festival of lamps, then the Sanskrit for lights, then say \emph{maṇḍayati} again
\end{itemize}

\begin{tcolorbox}[breakable,colback=teal!4,colframe=teal!35!black,title={Before you move on}]
What does \sk{मण्डयति} mean? (\textquotedblleft{}Decorates; adorns\textquotedblright{}.) Say it once more.
\end{tcolorbox}
\section[\texorpdfstring{\sk{कौतुकम्} (kautukam)}{कौतुकम् (kautukam)}]{\sk{कौतुकम्} (kautukam) — festive excitement; eagerness}

\label{lesson:SA-C326-kautukam}



\begin{quote}
Before the new one: say the Sanskrit for lights, then the Sanskrit for decorates.
\end{quote}

\subsection*{You'll want to know: \sk{कौतुकम्}}
\textbf{\sk{कौतुकम्}} — \emph{kautukam} — \textquotedblleft{}festive excitement; eagerness\textquotedblright{}.

\begin{quote}
\sk{उत्सवे} \sk{बालस्य} \sk{कौतुकं} \sk{वर्धते।}
\end{quote}

\begin{quote}
\emph{utsave bālasya kautukaṃ vardhate.}
\end{quote}

\begin{quote}
At the festival the child's excitement grows.
\end{quote}

A neuter noun.

\begin{cousinweb}
\textbf{\sk{कौतुकम्}} (\emph{kautukam}) is built on \emph{kutūhala}, \textquotedblleft{}curiosity\textquotedblright{} (MW abbreviates it), and MW gives \emph{kautuka} as curiosity, eagerness, and festivity, a show (MW).
\end{cousinweb}

\begin{grammarlens}[title={What you've built}]
One more word for telling what happened.
\end{grammarlens}

\subsection*{Your turn}
\begin{itemize}
\raggedright
  \item \emph{Say it:} \emph{kautukam}
  \item \emph{Say it:} \emph{kautukam}, once more
  \item \emph{Say it:} \emph{utsave bālasya kautukaṃ vardhate.} — At the festival the child's excitement grows.
  \item \emph{Recall:} say the Sanskrit for lights, then the Sanskrit for decorates, then say \emph{kautukam} again
\end{itemize}

\begin{tcolorbox}[breakable,colback=teal!4,colframe=teal!35!black,title={Before you move on}]
What does \sk{कौतुकम्} mean? (\textquotedblleft{}Festive excitement; eagerness\textquotedblright{}.) Say it once more.
\end{tcolorbox}
\section[\texorpdfstring{\sk{शोभा} (śobhā)}{शोभा (śobhā)}]{\sk{शोभा} (śobhā) — splendour; beauty}

\label{lesson:SA-C326-sobha}



\begin{quote}
Before the new one: say the Sanskrit for decorates, then the Sanskrit for festive excitement.
\end{quote}

\subsection*{You'll want to know: \sk{शोभा}}
\textbf{\sk{शोभा}} — \emph{śobhā} — \textquotedblleft{}splendour; beauty\textquotedblright{}.

\begin{quote}
\sk{दीपैः} \sk{नगरस्य} \sk{शोभा} \sk{वर्धते।}
\end{quote}

\begin{quote}
\emph{dīpaiḥ nagarasya śobhā vardhate.}
\end{quote}

\begin{quote}
With the lamps the splendour of the city grows.
\end{quote}

A feminine noun.

\begin{cousinweb}
\textbf{\sk{शोभा}} (\emph{śobhā}) is usually traced to the root √śubh, \textquotedblleft{}to shine\textquotedblright{}, as in \emph{śobhate}; MW gives \emph{śobhā} as splendour, brilliance, beauty, grace (MW).
\end{cousinweb}

\begin{grammarlens}[title={What you've built}]
One more word for telling what happened.
\end{grammarlens}

\subsection*{Your turn}
\begin{itemize}
\raggedright
  \item \emph{Say it:} \emph{śobhā}
  \item \emph{Say it:} \emph{śobhā}, once more
  \item \emph{Say it:} \emph{dīpaiḥ nagarasya śobhā vardhate.} — With the lamps the splendour of the city grows.
  \item \emph{Recall:} say the Sanskrit for decorates, then the Sanskrit for festive excitement, then say \emph{śobhā} again
\end{itemize}

\begin{tcolorbox}[breakable,colback=teal!4,colframe=teal!35!black,title={Before you move on}]
What does \sk{शोभा} mean? (\textquotedblleft{}Splendour; beauty\textquotedblright{}.) Say it once more.
\end{tcolorbox}
